\NativeSlot{47}{62}{%
% p054_s01

\ExamQuestion{10.}{5}%
若\(A\)表示1，\allowbreak{}\(B\)表示2，则\(C\)表示\mbox{（\hspace{1.8em}）}；若\(A\)表示2.3，\allowbreak{}\(B\)表示2.4，则\(D\)表示\mbox{（\hspace{1.8em}）}；若\(A\)表示48.6，\allowbreak{}\(B\)表示48.7，则\(C\)表示\mbox{（\hspace{1.8em}）}。\par
\BookFigure{figures/p054_s01-01.tex}
}
\NativeSlot{153}{61}{%
\ExamGroup{二、选择（将正确答案的序号填在括号里）。{\normalfont\kaishu\fontsize{10.8bp}{14bp}\selectfont （每空2分，共12分）}}

% p054_s02

\ExamQuestion{1.}{5}%
三星堆出土的青铜神树高3.96米，是全球最大单件青铜文物。3.96米中的“9”表示\mbox{（\hspace{1.8em}）}。\par
\ExamOptions{A. 9米}{B. 0.9分米}{C. 9分米}{D. 0.9厘米}
}
\NativeSlot{256.0}{71}{%
% p054_s03

\ExamQuestion{2.}{5}%
下列说法正确的是\mbox{（\hspace{1.8em}）}。\par
\ExamOptions{A. 用“四舍五入”法保留两位小数时，关键要看百分位上的数}{B. 一个两位小数精确到十分位约是3.5，这个数最小是3.50}{C. 一个小数的小数位数越多，这个小数就越大}{D. 所有的小数都有整数部分}
}
\NativeSlot{379}{58}{%
% p054_s04

\ExamQuestion{3.}{5}%
计算\(0.2\times0.03=(2\times0.1)\times(3\times0.01)=(2\times3)\times(0.1\times0.01)=6\times0.001=0.006\)，第三步的“6”表示的是新的\mbox{（\hspace{1.8em}）}，“0.001”表示的是新的\mbox{（\hspace{1.8em}）}。\par
\ExamOptions{A. 计数单位}{B. 计数单位的个数}{C. 积的小数位数}{D. 积}
}
\NativeSlot{477.0}{83}{%
% p054_s05

\ExamQuestion{4.}{5}%
学校运动会上，三名同学参加100米短跑比赛，成绩如下表。根据比赛排名，浩浩获得了一等奖，那么方框里最大能填\mbox{（\hspace{1.8em}）}。\par
\ExamOptions{A. 3}{B. 4}{C. 5}{D. 6}
\par\smallskip\Needspace{4\baselineskip}\begingroup\small\setlength{\tabcolsep}{3pt}
\noindent\begin{tabularx}{\linewidth}{|*{4}{>{\centering\arraybackslash}X|}}\hline
选手 & 明明 & 亮亮 & 浩浩\\\hline
成绩（秒） & 15.63 & 15.47 & \(15.\square6\)\\\hline
\end{tabularx}\par
\endgroup\smallskip
}
\NativeSlot{616.0}{36}{%
% p054_s06

\ExamQuestion{5.}{5}%
计算“8米6厘米＋560厘米”时，用“米”作单位，正确的算式是\mbox{（\hspace{1.8em}）}。\par
\ExamOptions{A. \(8.6+5.6\)}{B. \(8.06+5.6\)}{C. \(8.06+5.06\)}{D. \(8.6+5.06\)}
}
\NativeSlot{667.0}{72}{%
\ExamGroup{三、计算。{\normalfont\kaishu\fontsize{10.8bp}{14bp}\selectfont （共20分）}}

% p054_s07

\ExamQuestion{1.}{5}%
直接写得数。\par
\noindent \ExamOralCell{40\div100=}\ExamOralCell{0.6\times1000=}\ExamOralCell{3.6\div10=}\par
\noindent \ExamOralCell{45\div10\times100=}\ExamOralCell{0.21\div10=}\ExamOralCell{250\div1000=}\par
\noindent \ExamOralCell{4800\div100=}\ExamOralCell{3.2\times1000\div10=}\par
}
